% LaTeX's eqnarray as in the PDF: each row numbered, unless it has
% \nonumber, and a \ref to a row's \label gives the row's number.
% eqnarray* numbers no row, and a \label there gets the next number. A
% \qedhere in an eqnarray, which amsthm sets in the row, takes no number
% from the row. A \label on an unnumbered row gets the next row's number,
% also when the next row has a \label. At 11pt, the rows are \jot apart,
% which MathJax's default spacing is not.
\documentclass[11pt]{article}
\usepackage{amsmath,amsthm}
\begin{document}
\section{S}
\begin{eqnarray}
a & = & b \label{eq:first} \\
  & = & c \nonumber \\
  & = & d \label{eq:last}
\end{eqnarray}
\begin{eqnarray*}
e & = & f \label{eq:star} \\
  & = & g \label{eq:star2}
\end{eqnarray*}
Star \ref{eq:star} and \ref{eq:star2}.
\begin{proof}
\begin{eqnarray}
j & = & k \qedhere
\end{eqnarray}
\end{proof}
\begin{equation} h = i \label{eq:after} \end{equation}
See \ref{eq:first}, \ref{eq:last}, and \ref{eq:after}.
\begin{eqnarray}
l & = & m \nonumber \label{eq:skip} \\
  & = & n \label{eq:next}
\end{eqnarray}
Then \ref{eq:skip} and \ref{eq:next}.
\end{document}
